%%%%%
%
\section{FSSP: Activit� � partir l'enregistrement brut}

\itemstyle{Op\'erateur~:} \operfmt{fssp\_activ}
\index{fssp\_activ}

%%%%%OBJET%%%%
\itemstyle{Objet~:}
Calcul de l'activit� d'un FSSP exprim�e en proportion (gamme $[0, 1]$)
� partir de l'enregistrement brut de la voie "activit�".

L'enregistrement brut est bas� sur le comptage des pulsations
d'horloge dans les p�riodes pendant lesquelles la sonde
travaille (d�tection de gouttes, dimensionnement ou rejet).

Dans la sonde, le nombre de pulsations est divis�
par $2^x$, donc le coefficient est �gal �
$2^x$ / (fr�quence de l'horloge).


%%%%%ARGUMENTS%%%%
\bigskip
%%%%%ENTREES%%%%
\begin{tabarg}
\multicolumn{3}{l}{\titretab {Entr\'ees}}\\
\hline
$A_c$ & Vecteur, \texttt{float} & Activit\'e brute 
(\textit{comptes par p\'eriode d'\'echantillonnage})\\
$C_A$ & Coeff., \texttt{float} & Coefficient diviseur (1/ nombre de comptage max par seconde) \\
$f_e$ & Coeff., \texttt{float} & Fr\'equence d'\'echantillonnage ($Hz$)\\
\end{tabarg}

\medskip
%%%%%SORTIES%%%%
\begin{tabarg}
\multicolumn{3}{l}{\titretab {Sorties}}\\
\hline
$A_m$ & Vecteur, \texttt{float} & Activit\'e (proportion : $s$/$s$)\\
\end{tabarg}


%%%%%FORMULE%%%%
\itemstyle{Formule}

%%fssp_activ
\[
A_m = C_A.f_e.A_c
\]


%%%%%COMMENTAIRE%%%%
\itemstyle{Commentaire~:}
La valeur de $C_A$ pour le FSSP du CAM/SAFIRE est de 8.02e-6.

Pour d'autres FSSP, on trouve plut�t 1.00e-3.

%%%%%AUTEUR%%%%
\itemstyle{Auteur~:}
CNRM/GMEI/TRAMM

%%%%%REFERENCE%%%%
\itemstyle{R\'ef\'erences~:}
Brenguier-89 \cite{Brenguier-89}
